\begin{theorem}[Autocovariance of an MA(1) Process]
Let $\theta,\sigma^2\in\mathbb{R}$ with $\sigma^2\ge 0$. Let
$(\varepsilon_t)_{t\in\mathbb{Z}}$ be real-valued innovations with
\[
\mathbb{E}[\varepsilon_t]=0,
\qquad
\operatorname{Var}(\varepsilon_t)=\sigma^2
\]
for every $t$, and assume that distinct innovations are uncorrelated:
\[
\operatorname{Cov}(\varepsilon_s,\varepsilon_t)=0
\qquad
\text{whenever }s\ne t.
\]
Define the moving-average process
\[
X_t=\varepsilon_t+\theta\varepsilon_{t-1}.
\]
For $h\ge 0$, define
\[
\gamma(h)=\operatorname{Cov}(X_t,X_{t+h}).
\]
Then $\gamma(h)$ does not depend on $t$, and
\[
\gamma(0)=(1+\theta^2)\sigma^2,
\qquad
\gamma(1)=\theta\sigma^2,
\qquad
\gamma(h)=0\quad\text{for every }h\ge 2.
\]
\end{theorem}
